% Comparación trasladada detrás de la ley y de la ontología sectorial.
% Se excluyen únicamente la portada autónoma y el metalenguaje de edición.
\section{Evaluación cuantitativa posterior a la fijación del grafo generativo}
La tabla siguiente coloca en columnas contiguas toda coordenada escalar que la teoría permite evaluar de forma reproducible. La referencia experimental se abre después de fijar la salida HMT--MD. Una fila rotulada como evaluación condicionada afirma exactamente el valor de la función una vez fijada la firma; no atribuye a esa firma un selector físico que no haya sido construido. Los quarks se imprimen sin residual porque una comparación numérica requiere que esquema y escala de renormalización coincidan.

El censo numérico de esta tabla es exhaustivo respecto de las coordenadas escalares actualmente reproducibles: un electrón beta cerrado, diecisiete evaluaciones condicionadas ---once elementales y seis compuestas--- y dos salidas nulas exactas. Por tanto contiene veinte filas numéricas. Este número no es el dominio de la ley. Las clases cuya salida natural sigue siendo un operador, un espectro, una órbita o una multisección aparecen después con ese objeto completo, sin rellenar su columna con la referencia externa. El inventario experimental de 471 masas y 384 anchuras se conserva íntegro como contraste estado por estado; no se hace pasar su cardinal por un censo de predicciones HMT.

\Needspace{7\baselineskip}
\begingroup
\footnotesize
\setlength{\tabcolsep}{2.1pt}
\begin{longtable}{L{1.17cm}L{2.65cm}L{3.12cm}L{2.18cm}L{5.69cm}}
\toprule
\rowcolor{HMTNavy}\HMTtablehead{Estado} & \HMTtablehead{Valor HMT reproducible (MeV/\allowbreak{}c\({}^{2}\))} & \HMTtablehead{Referencia posterior} & \HMTtablehead{Diferencia relativa (ppm)} & \HMTtablehead{Naturaleza y alcance} \\
\midrule
\endfirsthead
\multicolumn{5}{l}{\sffamily\scriptsize\color{HMTGray}Continuación de la tabla}\\[-1pt]
\toprule
\rowcolor{HMTNavy}\HMTtablehead{Estado} & \HMTtablehead{Valor HMT reproducible (MeV/\allowbreak{}c\({}^{2}\))} & \HMTtablehead{Referencia posterior} & \HMTtablehead{Diferencia relativa (ppm)} & \HMTtablehead{Naturaleza y alcance} \\
\midrule
\endhead
\midrule
\endfoot
\bottomrule
\endlastfoot
e-/\allowbreak{}e+ & 0.510998\allowbreak{}950690\allowbreak{}000809 & 0.510998\allowbreak{}95069±\allowbreak{}0.000000\allowbreak{}00016 MeV/\allowbreak{}c\({}^{2}\) & 1.582406\allowbreak{}884×10\({}^{-9}\) & coordenada beta cerrada en la fibra central \\
\rowcolor{HMTLight}μ-/\allowbreak{}μ+ & 105.658360\allowbreak{}294435\allowbreak{}829303 & 105.658375\allowbreak{}5±\allowbreak{}0.000002\allowbreak{}3 MeV/\allowbreak{}c\({}^{2}\) & -0.143912\allowbreak{}53 & refinamiento E3 condicionado a la firma \\
τ-/\allowbreak{}τ+ & 1776.860917\allowbreak{}631432\allowbreak{}295595 & 1776.932465\allowbreak{}134090\allowbreak{}1 MeV/\allowbreak{}c\({}^{2}\) (valor central usado; presentación redondeada: 1776.93±\allowbreak{}0.09 MeV/\allowbreak{}c\({}^{2}\)) & -40.264615\allowbreak{}601 & refinamiento E3 condicionado a la firma \\
\rowcolor{HMTLight}u/\allowbreak{}\(\bar{\mathrm{u}}\) & 2.165924\allowbreak{}536173\allowbreak{}907663 & 2.16±\allowbreak{}0.07 MeV/\allowbreak{}c\({}^{2}\) & --- & carta angular; no se ha declarado firma Z\({}^{5}\) ni refinamiento E3 \\
d/\allowbreak{}\(\bar{\mathrm{d}}\) & 4.679550\allowbreak{}162948\allowbreak{}874172 & 4.70±\allowbreak{}0.07 MeV/\allowbreak{}c\({}^{2}\) & --- & carta angular; no se ha declarado firma Z\({}^{5}\) ni refinamiento E3 \\
\rowcolor{HMTLight}c/\allowbreak{}\(\bar{\mathrm{c}}\) & 1270.500838\allowbreak{}641911\allowbreak{}135286 & 1.2729±\allowbreak{}0.0045 GeV/\allowbreak{}c\({}^{2}\) & --- & carta angular; no se ha declarado firma Z\({}^{5}\) ni refinamiento E3 \\
s/\allowbreak{}\(\bar{\mathrm{s}}\) & 92.940093\allowbreak{}907845\allowbreak{}640641 & 92.9 ±\allowbreak{}0.7 MeV/\allowbreak{}c\({}^{2}\) & --- & carta angular; no se ha declarado firma Z\({}^{5}\) ni refinamiento E3 \\
\rowcolor{HMTLight}t/\allowbreak{}\(\bar{\mathrm{t}}\) & 172597.479544\allowbreak{}019136\allowbreak{}261367 & 172.60±\allowbreak{}0.27 GeV/\allowbreak{}c\({}^{2}\) & --- & carta angular; no se ha declarado firma Z\({}^{5}\) ni refinamiento E3 \\
b/\allowbreak{}\(\bar{\mathrm{b}}\) & 4190.119763\allowbreak{}299790\allowbreak{}218075 & 4.186±\allowbreak{}0.006 GeV/\allowbreak{}c\({}^{2}\) & --- & carta angular; no se ha declarado firma Z\({}^{5}\) ni refinamiento E3 \\
\rowcolor{HMTLight}γ & 0 & <1×10\({}^{-18}\) eV/\allowbreak{}c\({}^{2}\) & --- & salida nula exacta \\
g & 0 & --- & --- & salida nula exacta \\
\rowcolor{HMTLight}W+/\allowbreak{}W- & 80378.964909\allowbreak{}704547\allowbreak{}024865 & 80.3625 ±\allowbreak{}0.0077 GeV/\allowbreak{}c\({}^{2}\) & 204.882995\allowbreak{}235 & refinamiento E3 condicionado a la firma \\
Z0 & 91187.533444\allowbreak{}313407\allowbreak{}844855 & 91.187873\allowbreak{}297221\allowbreak{}344 GeV/\allowbreak{}c\({}^{2}\) (valor central usado; presentación redondeada: 91.1879±\allowbreak{}0.0020 GeV/\allowbreak{}c\({}^{2}\)) & -3.726952\allowbreak{}89 & refinamiento E3 condicionado a la firma \\
\rowcolor{HMTLight}H & 125292.466042\allowbreak{}536133\allowbreak{}000433 & 125.130943\allowbreak{}828161\allowbreak{}51 GeV/\allowbreak{}c\({}^{2}\) (valor central usado; presentación redondeada: 125.13±\allowbreak{}0.11 GeV/\allowbreak{}c\({}^{2}\)) & 1290.825509\allowbreak{}927 & carta angular; no se ha declarado firma Z\({}^{5}\) ni refinamiento E3 \\
π\({}^{0}\) & 134.976724\allowbreak{}327872\allowbreak{}125739 & 134.976827\allowbreak{}767685 MeV/\allowbreak{}c\({}^{2}\) & -0.766352\allowbreak{}378 & refinamiento E3 condicionado a la firma compuesta \\
\rowcolor{HMTLight}π±\allowbreak{} & 139.570485\allowbreak{}171572\allowbreak{}191375 & 139.570390\allowbreak{}983681 MeV/\allowbreak{}c\({}^{2}\) & 0.674841\allowbreak{}494 & refinamiento E3 condicionado a la firma compuesta \\
K±\allowbreak{} & 493.677085\allowbreak{}649530\allowbreak{}612476 & 493.676599\allowbreak{}458041 MeV/\allowbreak{}c\({}^{2}\) & 0.984838\allowbreak{}03 & refinamiento E3 condicionado a la firma compuesta \\
\rowcolor{HMTLight}K\({}^{0}\) & 497.611186\allowbreak{}497750\allowbreak{}457359 & 497.610767\allowbreak{}531258 MeV/\allowbreak{}c\({}^{2}\) & 0.841956\allowbreak{}243 & refinamiento E3 condicionado a la firma compuesta \\
protón & 938.271751\allowbreak{}546984\allowbreak{}898299 & 938.272089\allowbreak{}43 MeV/\allowbreak{}c\({}^{2}\) & -0.360111\allowbreak{}975 & refinamiento E3 condicionado a la firma compuesta \\
\rowcolor{HMTLight}neutrón & 939.565810\allowbreak{}472507\allowbreak{}023313 & 939.565421\allowbreak{}94 MeV/\allowbreak{}c\({}^{2}\) & 0.413523\allowbreak{}633 & refinamiento E3 condicionado a la firma compuesta \\
\end{longtable}
\endgroup
\Needspace{7\baselineskip}
\section{Evaluación individual de \(20\) secciones materiales y trazabilidad de sus coordenadas metrológicas}
\Needspace{7\baselineskip}
\subsection{Par electrón–positrón: ruta central, conjugación de hojas y carácter de masa}
\begin{itemize}
\item \textbf{Identificador catalográfico:} SM-LC-1.
\item \textbf{Tipo de descriptor:} selector basal y lectores bidireccionales.
\item \textbf{Descriptor:} selector (-22,+15); K\_D=\allowbreak{}K\_Ω=\allowbreak{}(80,54,6).
\item \textbf{Etapa:} beta bidireccional; coordenada inicial 0.510998\allowbreak{}922488\allowbreak{}066073 MeV/c\({}^{2}\).
\item \textbf{Evaluación HMT--MD reproducible:} 0.510998\allowbreak{}950690\allowbreak{}000809 MeV/c\({}^{2}\).
\item \textbf{Naturaleza matemática:} coordenada beta cerrada en la fibra central.
\item \textbf{Referencia posterior:} 0.510998\allowbreak{}95069±\allowbreak{}0.000000\allowbreak{}00016 MeV/c\({}^{2}\).
\item \textbf{Identidad y metrología externas:} Particle Data Group (PDG); PDG2026:\allowbreak{}observable:\allowbreak{}S003:\allowbreak{}S003M:\allowbreak{}332838; esquema no conservado en el corte documental; escala no conservada en el corte documental; nivel de confianza no conservado en el corte documental.
\item \textbf{Diferencia relativa:} 1.582406\allowbreak{}884×10\({}^{-9}\) ppm.
\item \textbf{Alcance:} derivación escalar cerrada.
\end{itemize}
\Needspace{7\baselineskip}
\subsection{Par muón–antimuón: conjugación de rutas y carácter leptónico}
\begin{itemize}
\item \textbf{Identificador catalográfico:} SM-LC-2.
\item \textbf{Tipo de descriptor:} firma completa σ\(\in\)Z\({}^{5}\).
\item \textbf{Descriptor:} (42,-3,8,0,2).
\item \textbf{Etapa:} torre beta E3; coordenada inicial 105.564059\allowbreak{}012859\allowbreak{}26472 MeV/c\({}^{2}\).
\item \textbf{Evaluación HMT--MD reproducible:} 105.658360\allowbreak{}294435\allowbreak{}829303 MeV/c\({}^{2}\).
\item \textbf{Naturaleza matemática:} refinamiento E3 condicionado a la firma.
\item \textbf{Referencia posterior:} 105.658375\allowbreak{}5±\allowbreak{}0.000002\allowbreak{}3 MeV/c\({}^{2}\).
\item \textbf{Identidad y metrología externas:} Particle Data Group (PDG); PDG2026:\allowbreak{}observable:\allowbreak{}S004:\allowbreak{}S004M:\allowbreak{}332849; esquema no conservado en el corte documental; escala no conservada en el corte documental; nivel de confianza no conservado en el corte documental.
\item \textbf{Diferencia relativa:} -0.143912\allowbreak{}53 ppm.
\item \textbf{Alcance:} evaluación exacta al fijar la firma; la selección física de esa firma constituye una operación independiente.
\end{itemize}
\Needspace{7\baselineskip}
\subsection{Par tau–antitau: conjugación de rutas y carácter leptónico}
\begin{itemize}
\item \textbf{Identificador catalográfico:} SM-LC-3.
\item \textbf{Tipo de descriptor:} firma completa σ\(\in\)Z\({}^{5}\).
\item \textbf{Descriptor:} (64,0,25,-1,6).
\item \textbf{Etapa:} torre beta E3; coordenada inicial 1771.945393\allowbreak{}886234\allowbreak{}110148 MeV/c\({}^{2}\).
\item \textbf{Evaluación HMT--MD reproducible:} 1776.860917\allowbreak{}631432\allowbreak{}295595 MeV/c\({}^{2}\).
\item \textbf{Naturaleza matemática:} refinamiento E3 condicionado a la firma.
\item \textbf{Referencia posterior:} 1776.932465\allowbreak{}134090\allowbreak{}1 MeV/c\({}^{2}\) (valor central usado; presentación redondeada: 1776.93±\allowbreak{}0.09 MeV/c\({}^{2}\)).
\item \textbf{Identidad y metrología externas:} Particle Data Group (PDG); PDG2026:\allowbreak{}observable:\allowbreak{}S035:\allowbreak{}S035M:\allowbreak{}332893; esquema no conservado en el corte documental; escala no conservada en el corte documental; nivel de confianza no conservado en el corte documental.
\item \textbf{Diferencia relativa:} -40.264615\allowbreak{}601 ppm.
\item \textbf{Alcance:} evaluación exacta al fijar la firma; la selección física de esa firma constituye una operación independiente.
\end{itemize}
\Needspace{7\baselineskip}
\subsection{Quark arriba y antiquark arriba}
\begin{itemize}
\item \textbf{Identificador catalográfico:} SM-Q-U1.
\item \textbf{Tipo de descriptor:} coindexación angular (n\_A,s\_C)\(\in\)Z\({}^{2}\).
\item \textbf{Descriptor:} (11,7).
\item \textbf{Etapa:} carta angular condicionada; coordenada inicial 2.165924\allowbreak{}536173\allowbreak{}907663 MeV/c\({}^{2}\).
\item \textbf{Evaluación HMT--MD reproducible:} 2.165924\allowbreak{}536173\allowbreak{}907663 MeV/c\({}^{2}\).
\item \textbf{Naturaleza matemática:} carta angular; no se ha declarado firma Z\({}^{5}\) ni refinamiento E3.
\item \textbf{Referencia posterior:} 2.16±\allowbreak{}0.07 MeV/c\({}^{2}\).
\item \textbf{Identidad y metrología externas:} Particle Data Group (PDG); PDG2026:\allowbreak{}observable:\allowbreak{}Q002:\allowbreak{}Q002M:\allowbreak{}333588; esquema no conservado en el corte documental; escala no conservada en el corte documental; nivel de confianza 90\%.
\item \textbf{Diferencia relativa:} --- ppm.
\item \textbf{Alcance:} evaluación condicionada; la firma no se ha elegido desde el valor comparado; esquema y escala externos requeridos para el residual.
\end{itemize}
\Needspace{7\baselineskip}
\subsection{Quark abajo y antiquark abajo}
\begin{itemize}
\item \textbf{Identificador catalográfico:} SM-Q-D1.
\item \textbf{Tipo de descriptor:} coindexación angular (n\_A,s\_C)\(\in\)Z\({}^{2}\).
\item \textbf{Descriptor:} (17,8).
\item \textbf{Etapa:} carta angular condicionada; coordenada inicial 4.679550\allowbreak{}162948\allowbreak{}874172 MeV/c\({}^{2}\).
\item \textbf{Evaluación HMT--MD reproducible:} 4.679550\allowbreak{}162948\allowbreak{}874172 MeV/c\({}^{2}\).
\item \textbf{Naturaleza matemática:} carta angular; no se ha declarado firma Z\({}^{5}\) ni refinamiento E3.
\item \textbf{Referencia posterior:} 4.70±\allowbreak{}0.07 MeV/c\({}^{2}\).
\item \textbf{Identidad y metrología externas:} Particle Data Group (PDG); PDG2026:\allowbreak{}observable:\allowbreak{}Q001:\allowbreak{}Q001M:\allowbreak{}333589; esquema no conservado en el corte documental; escala no conservada en el corte documental; nivel de confianza 90\%.
\item \textbf{Diferencia relativa:} --- ppm.
\item \textbf{Alcance:} evaluación condicionada; la firma no se ha elegido desde el valor comparado; esquema y escala externos requeridos para el residual.
\end{itemize}
\Needspace{7\baselineskip}
\subsection{Quark encanto y antiquark encanto}
\begin{itemize}
\item \textbf{Identificador catalográfico:} SM-Q-U2.
\item \textbf{Tipo de descriptor:} coindexación angular (n\_A,s\_C)\(\in\)Z\({}^{2}\).
\item \textbf{Descriptor:} (61,8).
\item \textbf{Etapa:} carta angular condicionada; coordenada inicial 1270.500838\allowbreak{}641911\allowbreak{}135286 MeV/c\({}^{2}\).
\item \textbf{Evaluación HMT--MD reproducible:} 1270.500838\allowbreak{}641911\allowbreak{}135286 MeV/c\({}^{2}\).
\item \textbf{Naturaleza matemática:} carta angular; no se ha declarado firma Z\({}^{5}\) ni refinamiento E3.
\item \textbf{Referencia posterior:} 1.2729±\allowbreak{}0.0045 GeV/c\({}^{2}\).
\item \textbf{Identidad y metrología externas:} Particle Data Group (PDG); PDG2026:\allowbreak{}observable:\allowbreak{}Q004:\allowbreak{}Q004M:\allowbreak{}333604; esquema no conservado en el corte documental; escala no conservada en el corte documental; nivel de confianza 90\%.
\item \textbf{Diferencia relativa:} --- ppm.
\item \textbf{Alcance:} evaluación condicionada; la firma no se ha elegido desde el valor comparado; esquema y escala externos requeridos para el residual.
\end{itemize}
\Needspace{7\baselineskip}
\subsection{Quark extraño y antiquark extraño}
\begin{itemize}
\item \textbf{Identificador catalográfico:} SM-Q-D2.
\item \textbf{Tipo de descriptor:} coindexación angular (n\_A,s\_C)\(\in\)Z\({}^{2}\).
\item \textbf{Descriptor:} (41,-3).
\item \textbf{Etapa:} carta angular condicionada; coordenada inicial 92.940093\allowbreak{}907845\allowbreak{}640641 MeV/c\({}^{2}\).
\item \textbf{Evaluación HMT--MD reproducible:} 92.940093\allowbreak{}907845\allowbreak{}640641 MeV/c\({}^{2}\).
\item \textbf{Naturaleza matemática:} carta angular; no se ha declarado firma Z\({}^{5}\) ni refinamiento E3.
\item \textbf{Referencia posterior:} 92.9 ±\allowbreak{}0.7 MeV/c\({}^{2}\).
\item \textbf{Identidad y metrología externas:} Particle Data Group (PDG); PDG2026:\allowbreak{}observable:\allowbreak{}Q003:\allowbreak{}Q003M:\allowbreak{}333590; esquema no conservado en el corte documental; escala no conservada en el corte documental; nivel de confianza 90\%.
\item \textbf{Diferencia relativa:} --- ppm.
\item \textbf{Alcance:} evaluación condicionada; la firma no se ha elegido desde el valor comparado; esquema y escala externos requeridos para el residual.
\end{itemize}
\Needspace{7\baselineskip}
\subsection{Quark cima y antiquark cima}
\begin{itemize}
\item \textbf{Identificador catalográfico:} SM-Q-U3.
\item \textbf{Tipo de descriptor:} coindexación angular (n\_A,s\_C)\(\in\)Z\({}^{2}\).
\item \textbf{Descriptor:} (100,-1).
\item \textbf{Etapa:} carta angular condicionada; coordenada inicial 172597.479544\allowbreak{}019136\allowbreak{}261367 MeV/c\({}^{2}\).
\item \textbf{Evaluación HMT--MD reproducible:} 172597.479544\allowbreak{}019136\allowbreak{}261367 MeV/c\({}^{2}\).
\item \textbf{Naturaleza matemática:} carta angular; no se ha declarado firma Z\({}^{5}\) ni refinamiento E3.
\item \textbf{Referencia posterior:} 172.60±\allowbreak{}0.27 GeV/c\({}^{2}\).
\item \textbf{Identidad y metrología externas:} Particle Data Group (PDG); PDG2026:\allowbreak{}observable:\allowbreak{}Q007:\allowbreak{}Q007TP:\allowbreak{}333614; esquema no conservado en el corte documental; escala no conservada en el corte documental; nivel de confianza no conservado en el corte documental.
\item \textbf{Diferencia relativa:} --- ppm.
\item \textbf{Alcance:} evaluación condicionada; la firma no se ha elegido desde el valor comparado; esquema y escala externos requeridos para el residual.
\end{itemize}
\Needspace{7\baselineskip}
\subsection{Quark fondo y antiquark fondo}
\begin{itemize}
\item \textbf{Identificador catalográfico:} SM-Q-D3.
\item \textbf{Tipo de descriptor:} coindexación angular (n\_A,s\_C)\(\in\)Z\({}^{2}\).
\item \textbf{Descriptor:} (71,-5).
\item \textbf{Etapa:} carta angular condicionada; coordenada inicial 4190.119763\allowbreak{}299790\allowbreak{}218075 MeV/c\({}^{2}\).
\item \textbf{Evaluación HMT--MD reproducible:} 4190.119763\allowbreak{}299790\allowbreak{}218075 MeV/c\({}^{2}\).
\item \textbf{Naturaleza matemática:} carta angular; no se ha declarado firma Z\({}^{5}\) ni refinamiento E3.
\item \textbf{Referencia posterior:} 4.186±\allowbreak{}0.006 GeV/c\({}^{2}\).
\item \textbf{Identidad y metrología externas:} Particle Data Group (PDG); PDG2026:\allowbreak{}observable:\allowbreak{}Q005:\allowbreak{}Q005M:\allowbreak{}333611; esquema no conservado en el corte documental; escala no conservada en el corte documental; nivel de confianza 90\%.
\item \textbf{Diferencia relativa:} --- ppm.
\item \textbf{Alcance:} evaluación condicionada; la firma no se ha elegido desde el valor comparado; esquema y escala externos requeridos para el residual.
\end{itemize}
\Needspace{7\baselineskip}
\subsection{Descriptor espectral reproducible del fotón: sector nulo, helicidad y carácter electromagnético}
\begin{itemize}
\item \textbf{Identificador catalográfico:} SM-GA-EM.
\item \textbf{Tipo de descriptor:} invariante sectorial nulo.
\item \textbf{Descriptor:} carta nula.
\item \textbf{Etapa:} carta nula; coordenada inicial 0 MeV/c\({}^{2}\).
\item \textbf{Evaluación HMT--MD reproducible:} 0 MeV/c\({}^{2}\).
\item \textbf{Naturaleza matemática:} salida nula exacta.
\item \textbf{Referencia posterior:} <1×10\({}^{-18}\) eV/c\({}^{2}\).
\item \textbf{Identidad y metrología externas:} Particle Data Group (PDG); PDG2026:\allowbreak{}observable:\allowbreak{}S000:\allowbreak{}S000M:\allowbreak{}332461; esquema no conservado en el corte documental; escala no conservada en el corte documental; nivel de confianza no conservado en el corte documental.
\item \textbf{Diferencia relativa:} --- ppm.
\item \textbf{Alcance:} exacto en el sector nulo; la cota externa no se interpreta como una masa positiva.
\item \textbf{Cota fotónica:} referencia externa ilustrativa no comparable; el corte identifica el observable pero no conserva el nivel de confianza.
\end{itemize}
\Needspace{7\baselineskip}
\subsection{Realización del gluón como sección de la representación adjunta de \(SU(3)\)}
\begin{itemize}
\item \textbf{Identificador catalográfico:} SM-GA-GL.
\item \textbf{Tipo de descriptor:} invariante sectorial nulo.
\item \textbf{Descriptor:} carta nula.
\item \textbf{Etapa:} carta nula; coordenada inicial 0 MeV/c\({}^{2}\).
\item \textbf{Evaluación HMT--MD reproducible:} 0 MeV/c\({}^{2}\).
\item \textbf{Naturaleza matemática:} salida nula exacta.
\item \textbf{Referencia posterior:} ---.
\item \textbf{Identidad y metrología externas:} fuente no consignada; observable no individualizado; esquema no consignado; escala no consignada; nivel de confianza no consignado.
\item \textbf{Diferencia relativa:} --- ppm.
\item \textbf{Alcance:} exacto en el sector nulo; la nulidad pertenece a la carta de calibre de color.
\end{itemize}
\Needspace{7\baselineskip}
\subsection{Bosones \texorpdfstring{\(W^{+}\)}{W+} y \texorpdfstring{\(W^{-}\)}{W−}}
\begin{itemize}
\item \textbf{Identificador catalográfico:} SM-GA-W.
\item \textbf{Tipo de descriptor:} firma completa σ\(\in\)Z\({}^{5}\).
\item \textbf{Descriptor:} (94,-1,0,-2,4).
\item \textbf{Etapa:} torre beta E3; coordenada inicial 80381.592550\allowbreak{}220666\allowbreak{}304144 MeV/c\({}^{2}\).
\item \textbf{Evaluación HMT--MD reproducible:} 80378.964909\allowbreak{}704547\allowbreak{}024865 MeV/c\({}^{2}\).
\item \textbf{Naturaleza matemática:} refinamiento E3 condicionado a la firma.
\item \textbf{Referencia posterior:} 80.3625 ±\allowbreak{}0.0077 GeV/c\({}^{2}\).
\item \textbf{Identidad y metrología externas:} Particle Data Group (PDG); PDG2026:\allowbreak{}observable:\allowbreak{}S043:\allowbreak{}S043M:\allowbreak{}332464; esquema no conservado en el corte documental; escala no conservada en el corte documental; nivel de confianza no conservado en el corte documental.
\item \textbf{Diferencia relativa:} 204.882995\allowbreak{}235 ppm.
\item \textbf{Alcance:} evaluación exacta al fijar la firma; la selección física de esa firma constituye una operación independiente.
\item \textbf{Convención resonante:} el corte externo no fija conjuntamente convención de polo, Breit--Wigner o sobre capa, anchura y covarianza; el residual se interpreta sólo como información escalar.
\end{itemize}
\Needspace{7\baselineskip}
\subsection{Descriptor espectral reproducible del bosón \(Z^0\): mezcla neutra, polarización y carácter electrodébil}
\begin{itemize}
\item \textbf{Identificador catalográfico:} SM-GA-Z.
\item \textbf{Tipo de descriptor:} firma completa σ\(\in\)Z\({}^{5}\).
\item \textbf{Descriptor:} (95,-1,-11,0,-4).
\item \textbf{Etapa:} torre beta E3; coordenada inicial 91299.748286\allowbreak{}598170\allowbreak{}117317 MeV/c\({}^{2}\).
\item \textbf{Evaluación HMT--MD reproducible:} 91187.533444\allowbreak{}313407\allowbreak{}844855 MeV/c\({}^{2}\).
\item \textbf{Naturaleza matemática:} refinamiento E3 condicionado a la firma.
\item \textbf{Referencia posterior:} 91.187873\allowbreak{}297221\allowbreak{}344 GeV/c\({}^{2}\) (valor central usado; presentación redondeada: 91.1879±\allowbreak{}0.0020 GeV/c\({}^{2}\)).
\item \textbf{Identidad y metrología externas:} Particle Data Group (PDG); PDG2026:\allowbreak{}observable:\allowbreak{}S044:\allowbreak{}S044M:\allowbreak{}332513; esquema no conservado en el corte documental; escala no conservada en el corte documental; nivel de confianza no conservado en el corte documental.
\item \textbf{Diferencia relativa:} -3.726952\allowbreak{}89 ppm.
\item \textbf{Alcance:} evaluación exacta al fijar la firma; la selección física de esa firma constituye una operación independiente.
\item \textbf{Convención resonante:} el corte externo no fija conjuntamente convención de polo, Breit--Wigner o sobre capa, anchura y covarianza; el residual se interpreta sólo como información escalar.
\end{itemize}
\Needspace{7\baselineskip}
\subsection{Realización del bosón de Higgs: modo escalar, acoplamiento y carácter de masa}
\begin{itemize}
\item \textbf{Identificador catalográfico:} SM-H-1.
\item \textbf{Tipo de descriptor:} coindexación angular (n\_A,s\_C)\(\in\)Z\({}^{2}\).
\item \textbf{Descriptor:} (97,9).
\item \textbf{Etapa:} carta angular condicionada; coordenada inicial 125292.466042\allowbreak{}536133\allowbreak{}000433 MeV/c\({}^{2}\).
\item \textbf{Evaluación HMT--MD reproducible:} 125292.466042\allowbreak{}536133\allowbreak{}000433 MeV/c\({}^{2}\).
\item \textbf{Naturaleza matemática:} carta angular; no se ha declarado firma Z\({}^{5}\) ni refinamiento E3.
\item \textbf{Referencia posterior:} 125.130943\allowbreak{}828161\allowbreak{}51 GeV/c\({}^{2}\) (valor central usado; presentación redondeada: 125.13±\allowbreak{}0.11 GeV/c\({}^{2}\)).
\item \textbf{Identidad y metrología externas:} Particle Data Group (PDG); PDG2026:\allowbreak{}observable:\allowbreak{}S126:\allowbreak{}S126M:\allowbreak{}332733; esquema no conservado en el corte documental; escala no conservada en el corte documental; nivel de confianza no conservado en el corte documental.
\item \textbf{Diferencia relativa:} 1290.825509\allowbreak{}927 ppm.
\item \textbf{Alcance:} evaluación condicionada; la firma no se ha elegido desde el valor comparado.
\item \textbf{Convención resonante:} el corte externo no fija conjuntamente convención de polo, Breit--Wigner o sobre capa, anchura y covarianza; el residual se interpreta sólo como información escalar.
\end{itemize}
\Needspace{7\baselineskip}
\subsection{Pión neutro \texorpdfstring{\(\pi^{0}\)}{pi0}}
\begin{itemize}
\item \textbf{Identificador catalográfico:} CMP-PI0.
\item \textbf{Tipo de descriptor:} firma completa σ\(\in\)Z\({}^{5}\).
\item \textbf{Descriptor:} (44,-4,-25,1,-2).
\item \textbf{Etapa:} torre beta E3 compuesta; coordenada inicial 135.350257\allowbreak{}251504\allowbreak{}283349 MeV/c\({}^{2}\).
\item \textbf{Evaluación HMT--MD reproducible:} 134.976724\allowbreak{}327872\allowbreak{}125739 MeV/c\({}^{2}\).
\item \textbf{Naturaleza matemática:} refinamiento E3 condicionado a la firma compuesta.
\item \textbf{Referencia posterior:} 134.976827\allowbreak{}767685 MeV/c\({}^{2}\).
\item \textbf{Identidad y metrología externas:} referencia compuesta conservada por el suplemento E3; la convención, la covarianza y la edición se consignan en la fuente metrológica de referencia.
\item \textbf{Diferencia relativa:} -0.766352\allowbreak{}378 ppm.
\item \textbf{Alcance:} evaluación exacta al fijar la firma compuesta; el generador de ligadura permanece como morfismo propio.
\end{itemize}
\Needspace{7\baselineskip}
\subsection{Piones cargados \(\pi^\pm\): composición quark–antiquark, carga y masa}
\begin{itemize}
\item \textbf{Identificador catalográfico:} CMP-PIPM.
\item \textbf{Tipo de descriptor:} firma completa σ\(\in\)Z\({}^{5}\).
\item \textbf{Descriptor:} (44,1,-2,2,-1).
\item \textbf{Etapa:} torre beta E3 compuesta; coordenada inicial 139.596265\allowbreak{}529971\allowbreak{}781015 MeV/c\({}^{2}\).
\item \textbf{Evaluación HMT--MD reproducible:} 139.570485\allowbreak{}171572\allowbreak{}191375 MeV/c\({}^{2}\).
\item \textbf{Naturaleza matemática:} refinamiento E3 condicionado a la firma compuesta.
\item \textbf{Referencia posterior:} 139.570390\allowbreak{}983681 MeV/c\({}^{2}\).
\item \textbf{Identidad y metrología externas:} referencia compuesta conservada por el suplemento E3; la convención, la covarianza y la edición se consignan en la fuente metrológica de referencia.
\item \textbf{Diferencia relativa:} 0.674841\allowbreak{}494 ppm.
\item \textbf{Alcance:} evaluación exacta al fijar la firma compuesta; el generador de ligadura permanece como morfismo propio.
\end{itemize}
\Needspace{7\baselineskip}
\subsection{Kaones cargados \(K^\pm\): composición con extrañeza, carga y masa}
\begin{itemize}
\item \textbf{Identificador catalográfico:} CMP-KPM.
\item \textbf{Tipo de descriptor:} firma completa σ\(\in\)Z\({}^{5}\).
\item \textbf{Descriptor:} (54,-1,17,-2,2).
\item \textbf{Etapa:} torre beta E3 compuesta; coordenada inicial 492.762503\allowbreak{}210140\allowbreak{}547854 MeV/c\({}^{2}\).
\item \textbf{Evaluación HMT--MD reproducible:} 493.677085\allowbreak{}649530\allowbreak{}612476 MeV/c\({}^{2}\).
\item \textbf{Naturaleza matemática:} refinamiento E3 condicionado a la firma compuesta.
\item \textbf{Referencia posterior:} 493.676599\allowbreak{}458041 MeV/c\({}^{2}\).
\item \textbf{Identidad y metrología externas:} referencia compuesta conservada por el suplemento E3; la convención, la covarianza y la edición se consignan en la fuente metrológica de referencia.
\item \textbf{Diferencia relativa:} 0.984838\allowbreak{}03 ppm.
\item \textbf{Alcance:} evaluación exacta al fijar la firma compuesta; el generador de ligadura permanece como morfismo propio.
\end{itemize}
\Needspace{7\baselineskip}
\subsection{Kaón neutro \texorpdfstring{\(K^{0}\)}{K0}}
\begin{itemize}
\item \textbf{Identificador catalográfico:} CMP-K0.
\item \textbf{Tipo de descriptor:} firma completa σ\(\in\)Z\({}^{5}\).
\item \textbf{Descriptor:} (54,0,32,3,-2).
\item \textbf{Etapa:} torre beta E3 compuesta; coordenada inicial 495.816066\allowbreak{}594426\allowbreak{}768595 MeV/c\({}^{2}\).
\item \textbf{Evaluación HMT--MD reproducible:} 497.611186\allowbreak{}497750\allowbreak{}457359 MeV/c\({}^{2}\).
\item \textbf{Naturaleza matemática:} refinamiento E3 condicionado a la firma compuesta.
\item \textbf{Referencia posterior:} 497.610767\allowbreak{}531258 MeV/c\({}^{2}\).
\item \textbf{Identidad y metrología externas:} referencia compuesta conservada por el suplemento E3; la convención, la covarianza y la edición se consignan en la fuente metrológica de referencia.
\item \textbf{Diferencia relativa:} 0.841956\allowbreak{}243 ppm.
\item \textbf{Alcance:} evaluación exacta al fijar la firma compuesta; el generador de ligadura permanece como morfismo propio.
\end{itemize}
\Needspace{7\baselineskip}
\subsection{Protón: composición bariónica, neutralización de color y carácter de masa}
\begin{itemize}
\item \textbf{Identificador catalográfico:} CMP-P.
\item \textbf{Tipo de descriptor:} firma completa σ\(\in\)Z\({}^{5}\).
\item \textbf{Descriptor:} (59,0,9,1,0).
\item \textbf{Etapa:} torre beta E3 compuesta; coordenada inicial 937.314779\allowbreak{}258699\allowbreak{}634563 MeV/c\({}^{2}\).
\item \textbf{Evaluación HMT--MD reproducible:} 938.271751\allowbreak{}546984\allowbreak{}898299 MeV/c\({}^{2}\).
\item \textbf{Naturaleza matemática:} refinamiento E3 condicionado a la firma compuesta.
\item \textbf{Referencia posterior:} 938.272089\allowbreak{}43 MeV/c\({}^{2}\).
\item \textbf{Identidad y metrología externas:} referencia compuesta conservada por el suplemento E3; la convención, la covarianza y la edición se consignan en la fuente metrológica de referencia.
\item \textbf{Diferencia relativa:} -0.360111\allowbreak{}975 ppm.
\item \textbf{Alcance:} evaluación exacta al fijar la firma compuesta; el generador de ligadura permanece como morfismo propio.
\end{itemize}
\Needspace{7\baselineskip}
\subsection{Neutrón: composición bariónica, neutralización de color y carácter de masa}
\begin{itemize}
\item \textbf{Identificador catalográfico:} CMP-N.
\item \textbf{Tipo de descriptor:} firma completa σ\(\in\)Z\({}^{5}\).
\item \textbf{Descriptor:} (59,1,-34,0,1).
\item \textbf{Etapa:} torre beta E3 compuesta; coordenada inicial 943.123155\allowbreak{}648642\allowbreak{}014982 MeV/c\({}^{2}\).
\item \textbf{Evaluación HMT--MD reproducible:} 939.565810\allowbreak{}472507\allowbreak{}023313 MeV/c\({}^{2}\).
\item \textbf{Naturaleza matemática:} refinamiento E3 condicionado a la firma compuesta.
\item \textbf{Referencia posterior:} 939.565421\allowbreak{}94 MeV/c\({}^{2}\).
\item \textbf{Identidad y metrología externas:} referencia compuesta conservada por el suplemento E3; la convención, la covarianza y la edición se consignan en la fuente metrológica de referencia.
\item \textbf{Diferencia relativa:} 0.413523\allowbreak{}633 ppm.
\item \textbf{Alcance:} evaluación exacta al fijar la firma compuesta; el generador de ligadura permanece como morfismo propio.
\end{itemize}
% Las tablas históricas abreviadas permanecen conservadas a continuación en
% este propietario, pero quedan fuera del tronco científico vigente: no
% definen el dominio 81/56/13/324, la firma, el carácter ni las torres.
\endinput
\clearpage
\begin{landscape}
\section{Malla histórica extendida conservada}
\begingroup
\fontsize{6.8}{7.7}\selectfont
\setlength{\tabcolsep}{1.15pt}
\renewcommand{\arraystretch}{1.06}
\begin{longtable}{L{2.93cm}L{2.93cm}L{2.93cm}L{2.93cm}L{2.93cm}L{2.93cm}L{2.93cm}L{2.93cm}}
\toprule
\rowcolor{HMTNavy}\HMTtablehead{Estado} & \HMTtablehead{n} & \HMTtablehead{k} & \HMTtablehead{t} & \HMTtablehead{Evaluación HMT histórica (MeV/\allowbreak{}c\({}^{2}\))} & \HMTtablehead{Referencia posterior (MeV/\allowbreak{}c\({}^{2}\))} & \HMTtablehead{Diferencia consignada (ppm)} & \HMTtablehead{Fuerza matemática} \\
\midrule
\endfirsthead
\multicolumn{8}{l}{\sffamily\scriptsize\color{HMTGray}Continuación de la tabla}\\[-1pt]
\toprule
\rowcolor{HMTNavy}\HMTtablehead{Estado} & \HMTtablehead{n} & \HMTtablehead{k} & \HMTtablehead{t} & \HMTtablehead{Evaluación HMT histórica (MeV/\allowbreak{}c\({}^{2}\))} & \HMTtablehead{Referencia posterior (MeV/\allowbreak{}c\({}^{2}\))} & \HMTtablehead{Diferencia consignada (ppm)} & \HMTtablehead{Fuerza matemática} \\
\midrule
\endhead
\midrule
\endfoot
\bottomrule
\endlastfoot
mu & 64 & 146 & 1 & 105.656529 & 105.658374 & -17.461938\allowbreak{}2274 & evidencia numérica histórica condicionada; no constituye una regla física de selección anterior al contraste \\
\rowcolor{HMTLight}tau & 64 & 349 & -1 & 1776.835 & 1776.86 & -14.069763\allowbreak{}5154 & evidencia numérica histórica condicionada; no constituye una regla física de selección anterior al contraste \\
π±\allowbreak{} & 37 & 125 & 0 & 139.570582 & 139.57039 & 1.375649\allowbreak{}94982 & evidencia numérica histórica condicionada; no constituye una regla física de selección anterior al contraste \\
\rowcolor{HMTLight}K±\allowbreak{} & 52 & 177 & 1 & 493.679 & 493.677 & 4.051231\allowbreak{}87833 & evidencia numérica histórica condicionada; no constituye una regla física de selección anterior al contraste \\
K\({}^{0}\) & 53 & 178 & 1 & 497.616 & 497.611 & 10.048009\allowbreak{}3889 & evidencia numérica histórica condicionada; no constituye una regla física de selección anterior al contraste \\
\rowcolor{HMTLight}p & 55 & 195 & 1 & 938.27244 & 938.272081 & 0.382618\allowbreak{}226919 & evidencia numérica histórica condicionada; no constituye una regla física de selección anterior al contraste \\
n & 55 & 197 & 0 & 939.538 & 939.565413 & -29.176254\allowbreak{}9161 & evidencia numérica histórica condicionada; no constituye una regla física de selección anterior al contraste \\
\rowcolor{HMTLight}eta & 46 & 158 & 0 & 547.836 & 547.862 & -47.457206\allowbreak{}3768 & evidencia numérica histórica condicionada; no constituye una regla física de selección anterior al contraste \\
ρ\({}^{0}\) & 50 & 176 & 0 & 775.267 & 775.26 & 9.029228\allowbreak{}90385 & evidencia numérica histórica condicionada; no constituye una regla física de selección anterior al contraste \\
\rowcolor{HMTLight}φ(1020) & 54 & 190 & 0 & 1019.475 & 1019.461 & 13.732747\allowbreak{}0104 & evidencia numérica histórica condicionada; no constituye una regla física de selección anterior al contraste \\
J/\allowbreak{}ψ & 64 & 237 & 0 & 3096.979 & 3096.916 & 20.342818\allowbreak{}4684 & evidencia numérica histórica condicionada; no constituye una regla física de selección anterior al contraste \\
\rowcolor{HMTLight}Υ(1S) & 89 & 334 & 0 & 9460.34 & 9460.3 & 4.228195\allowbreak{}72318 & evidencia numérica histórica condicionada; no constituye una regla física de selección anterior al contraste \\
\end{longtable}
\endgroup
\end{landscape}
Estas doce cifras no se fusionan con las veinte salidas de la tabla rectora. La malla conserva cinco etiquetas adicionales ---η, ρ\({}^{0}\), φ(1020), J/ψ y Υ(1S)---, pero aún no construye el mapa estado compuesto→ruta y registro genealógico→(n,k,t). El cuanto interno CT108, 0.058177\allowbreak{}641733\allowbreak{}144319\allowbreak{}230789\allowbreak{}692282\allowbreak{}953757 en la normalización hbar\_int=\allowbreak{}c\_int=\allowbreak{}t0=\allowbreak{}1, es una sección de la recta de masa y no una especie.

La tabla no delimita el dominio. El dominio interno continúa siendo el atlas de 81 celdas, 56 familias, 13 multisecciones y 324 rutas orientadas; las realizaciones físicas tipadas y el inventario de 471 masas y 384 anchuras se desarrollan después.

\Needspace{7\baselineskip}
